\documentclass[12pt]{article}
\usepackage[utf8]{inputenc}
\usepackage{geometry}
\usepackage{booktabs}
\usepackage{graphicx}
\usepackage{hyperref}

\title{\textbf{Quantifying Layer 2 Convergence: A Comparative Performance Analysis of PVST+ and Rapid PVST+ in Redundant Cisco Topologies}}

\author{
    \textbf{Lead Author/Editor:} [Name] \\
    \textbf{Co-Authors:} [Names] \\
    \textit{Network Engineering Research Group}
}

\date{March 2026}

\begin{document}

\maketitle

\begin{abstract}
This paper investigates the performance differences between legacy Per-VLAN Spanning Tree Plus (PVST+) and Rapid PVST+ (802.1w) within a physical Cisco switching environment. While theoretical models suggest a significant reduction in convergence time, this study seeks to quantify that downtime through empirical packet-loss data. By simulating link failures in a redundant triangle topology, we analyze the impact of protocol tuning on network availability and professional infrastructure reliability.
\end{abstract}

\section{Introduction}
Modern local area networks (LANs) require high availability, often achieved through physical link redundancy. However, without a loop-prevention mechanism, redundant paths result in broadcast storms. 



The Spanning Tree Protocol (STP) prevents these loops by strategically blocking ports. This research critiques previous amateur-level simulations and provides high-fidelity data gathered from physical hardware to determine if the transition to Rapid-STP is mandatory for modern enterprise standards.

\section{Methodology}
The experimental setup consisted of three Cisco switches wired in a redundant triangle topology. 



\subsection{Environment Setup}
\begin{itemize}
    \item \textbf{Hardware:} 3x Cisco Catalyst Switches, 2x Linux End-Hosts.
    \item \textbf{Metric:} ICMP Echo Request/Reply loss using $ping -O$.
    \item \textbf{Variables:} We tested convergence under three conditions: Default PVST+, Rapid PVST+, and Optimized Rapid PVST+ with manual Root Bridge priority.
\end{itemize}

\section{Data and Results}
The following table represents the raw data collected during the "Link Cut" simulations.

\begin{center}
\begin{tabular}{lccc}
\toprule
\textbf{Protocol} & \textbf{Packets Sent} & \textbf{Packets Lost} & \textbf{Downtime (s)} \\ \midrule
PVST+ (Legacy)    & 100 & [X] & $\approx 50s$ \\
Rapid PVST+       & 100 & [Y] & $\approx 2s$ \\
Optimized R-PVST+ & 100 & [Z] & $< 1s$ \\ \bottomrule
\end{tabular}
\end{center}

\section{Analysis}


Conversely, the Rapid PVST+ handshake mechanism allows for nearly instantaneous transition by utilizing the \textit{Proposal/Agreement} process on point-to-point links.

\section{Conclusion}
Our findings confirm that legacy STP is insufficient for modern real-time traffic (VoIP, Video). Furthermore, we proved that manual Root Bridge selection is critical to prevent non-deterministic traffic paths. This activity serves as a baseline for future research into Link Aggregation (EtherChannel) and Layer 3 redundancy.

\section*{Author Contributions \& Acknowledgments}
\begin{itemize}
    \item \textbf{Co-Authors:} [Names] - Contributed to data analysis, technical writing, and critique.
    \item \textbf{Contributors:} [Names] - Conducted physical lab setup, wiring, and data collection.
    \item \textbf{Editors:} [Names] - Final technical review and formatting.
    \item \textbf{Acknowledgments:} We thank the club members who provided feedback during the Phase 1 "Stop \& Think" sessions.
\end{itemize}

\end{document}